\tabref{views}に，ImageNet-1Kのベンチマークで6つの特徴抽出器をそれぞれ単一のビューとして使った値を示す（抽出1通り）．DINOv2 B/14・L/14，DINOv3 B/16・L/16，DINO (v1) ViT-B/16，CLIP ViT-B/16の\NViewsN{}ビューのすべてで，$\sigma$はnear-OODの\FPR{}が静的$p$（\NViewsNMainBelowStaticNear{}/\NViewsN{}），LGC-max（\NViewsNMainBelowLgcNear{}/\NViewsN{}），$\pP$だけ（\NViewsNMainBelowLpNear{}/\NViewsN{}）より低く，交差適合なしの$\pP\pN$との差は\NViewsMinMainMinusVOneNear{}--+\NViewsMaxMainMinusVOneNear{}ポイントである．たとえばCLIP ViT-B/16（ベース検出器と同じ重み）では静的$p$ \NViewClipStaticNearFpr\%から$\sigma$ \NViewClipMainNearFpr\%へ，DINO (v1)では\NViewDinoOneStaticNearFpr\%から\NViewDinoOneMainNearFpr\%へ下がる．利得は自己教師あり学習の世代やVLMかどうかに依らず，特徴空間でIDが支持画像のまわりに集まり未知が互いに近いという性質だけに依る．\textbf{CLIPの静的$p$について．}CLIP ViT-B/16の静的$p$とkNNはnear \NViewClipStaticNearFpr\%，far \NViewClipStaticFarFpr\%とほとんど分離しない．特徴はTINSが使うのと同じ画像埋め込み（投影後，単位ノルム）であり取り違えではなく，同じ特徴でゼロショットのMCMはImageNet-200で動く（\tabref{in200}）．CLIPの画像埋め込みは狭い錐に集まり（支持画像の平均のノルムは\NViewsSupportMeanNormClip{}，DINO系では\NViewsSupportMeanNormDinoMin{}--\NViewsSupportMeanNormDinoMax{}），最も近い原型への余弦はIDで\NViewClipMaxCosIdMean{}$\pm$\NViewClipMaxCosIdStd{}，Texturesで\NViewClipMaxCosOodMean{}$\pm$\NViewClipMaxCosOodStd{}と平均は分かれるが，IDの下側の裾がTexturesの全体を覆うので，IDの95\%を受け入れる閾値ではTexturesのほぼ全部が通る．登録時の仮説「支持画像の平均を引いてから余弦を取れば（中心化）DINO系の範囲に入る」は否定された（付録\ref{sec:app_hyp}）．中心化で静的$p$はnear \NViewClipStaticcNearFpr\%，far \NViewClipStaticcFarFpr\%，kNNはnear \NViewClipKnncNearFpr\%，far \NViewClipKnncFarFpr\%になり，farは\NViewClipCentredMinusRawFarFprAbs{}ポイント改善するがnearはほぼ変わらない（\tabref{views}の中心化の列；DINO系の静的$p$の変化は\NViewsDinoCentredMaxabsFpr{}ポイント以内）．中心化なしの同じ計算（交差適合なし）は表の静的$p$と\NViewsRawVsTableMaxabs{}ポイント以内で一致する．グラフは相対的な近傍だけを使うので，同じ特徴でも$\pP$と$\sigma$はこの影響を受けない．
